% SPDX-License-Identifier: CC-BY-4.0
% Route B companion paper, Draft 2 (November 12, 2026).
% See OUTLINE.md for the plan and revision_log.md for the history.
\documentclass[12pt,letterpaper]{article}
\usepackage[T1]{fontenc}
% Times for text and math. The style guide prefers newtx. It was not
% installed where this draft was built, so the fallback mathptmx is used.
% To switch, replace the next line by \usepackage{newtxtext,newtxmath}.
\usepackage{mathptmx}
\usepackage[letterpaper,margin=1in]{geometry}
\usepackage{amsmath,amssymb,amsthm,mathrsfs}
\usepackage{enumitem,booktabs,longtable,array}
\usepackage{xcolor}
\usepackage[colorlinks=true,linkcolor=blue!50!black,citecolor=blue!50!black,urlcolor=blue!50!black]{hyperref}
\setlist{itemsep=2pt,topsep=4pt}
\emergencystretch=3em
\numberwithin{equation}{section}

\theoremstyle{plain}
\newtheorem{theorem}{Theorem}[section]
\newtheorem{proposition}[theorem]{Proposition}
\newtheorem{lemma}[theorem]{Lemma}
\newtheorem{corollary}[theorem]{Corollary}
\newtheorem{conjecture}[theorem]{Conjecture}
% Theorems A-D of the Introduction: a lettered counter, so that they can
% be cited with \ref.
\newtheorem{introtheorem}{Theorem}
\renewcommand{\theintrotheorem}{\Alph{introtheorem}}
\theoremstyle{definition}
\newtheorem{definition}[theorem]{Definition}
\newtheorem{example}[theorem]{Example}
\theoremstyle{remark}
\newtheorem{remark}[theorem]{Remark}
\newtheorem{status}[theorem]{Status}

\newcommand{\OO}{\mathcal{O}}
\newcommand{\Hc}{\mathcal{H}}
\newcommand{\cP}{\mathcal{P}}
\newcommand{\cL}{\mathcal{L}}
\newcommand{\fa}{\mathfrak{a}}
\newcommand{\ip}[1]{\langle #1\rangle}
\newcommand{\Th}{\widehat{T}}
\newcommand{\Thu}{\widehat{T}^{\,u}}
\DeclareMathOperator{\vol}{vol}
\DeclareMathOperator{\tr}{tr}
\DeclareMathOperator{\diag}{diag}
\DeclareMathOperator{\Ad}{Ad}
\newcommand{\one}{\mathbf 1}
\newcommand{\ellr}{\mathrm{ell}}
\newcommand{\Tell}{T_{\mathrm{ell}}}
\newcommand{\Tdisc}{T_{\mathrm{disc}}}

\title{The spectral measure of an elliptic orbital integral,\\ via Arthur's local trace formula\thanks{This is a working draft. Statements are labeled as proved here, quoted (with a reference), or expected; every ``Status'' paragraph says which, and what remains to be checked.}\\[1ex]
\large Draft 2}
\author{David Joyner\thanks{\texttt{wdjoyner@gmail.com}} \and
Claude\thanks{Claude is an AI model made by Anthropic. See the Acknowledgments for the division of work.}}
\date{First draft: November 10, 2026. Draft 2: November 12, 2026}

\begin{document}
\maketitle

\begin{abstract}
This paper was written with the help of an AI system: Claude (Anthropic) suggested the approach and drafted the paper as a co-author. See the Acknowledgments for details. The paper revises and updates the first author's earlier paper \cite{Jo92}, using a method different from the one used in the analogous paper \cite{JC26}. Let $G$ be an unramified semisimple group over a $p$-adic field, $K$ a hyperspecial maximal compact subgroup, and $\gamma$ a regular elliptic element. Taking one test function spherical and the other supported on the elliptic set, we specialize Arthur's local trace formula. On the geometric side only the orbital integrals of $G$ survive. On the spectral side only two kinds of terms survive: the unramified principal series of $G$, weighted by a weighted character of the elliptic test function, and Arthur's elliptic tempered virtual characters with $K$-fixed vectors. We deduce that the spherical orbital integrals of $\gamma$ are given by a finite measure on the unitary unramified characters. Its absolutely continuous part has density a weighted character. Its singular part is supported on subtori where the $R$-groups of the unramified principal series are non-trivial, and it vanishes when the center of $G$ is connected. For $SL(n)$ the singular part is given by $\kappa$-orbital integrals. We compare the result with the truncation method of the companion paper \cite{JC26}, and work out $PGL(2)$, $SL(2)$, $PGL(3)$ and $SL(3)$.

This is a working draft. The results are proved modulo normalizations in Arthur's formula that have not yet been checked. The main theorem is proved in outline, and the statement for $SL(n)$ is given in outline only. The section ``Status and what remains'' lists what is missing.
\end{abstract}

\tableofcontents

% =====================================================================
\section{Introduction}\label{s:intro}

This introduction states the problem and recalls the truncation method of the companion paper. It then describes the method of this paper, states the results, and gives a plan of the paper.

\paragraph{The problem.} Let $F$ be a $p$-adic field with ring of integers $\OO$ and residue field of order $q$. Let $G$ be a connected unramified semisimple group over $F$, and $K\subset G(F)$ a hyperspecial maximal compact subgroup. Let $\Thu$ be the compact torus of unitary unramified characters of a minimal Levi subgroup $M_0$, with Weyl group $W$. For $f$ in the spherical Hecke algebra $\Hc(G,K)$ let $f^\vee$ be its Satake transform, a $W$-invariant function on $\Thu$, and let $d\mu$ be the Plancherel measure on $\Thu/W$. For $\gamma\in G(F)$ regular elliptic, the orbital integral $f\mapsto O_\gamma(f)=\int_{G(F)}f(g^{-1}\gamma g)\,dg$ is a linear form on $\Hc(G,K)$, i.e.\ on $W$-invariant Laurent polynomials on $\Thu$. We ask what kind of object it is:
\begin{equation}\label{e:main}
O_\gamma(f)=\int_{\Thu/W}f^\vee(\chi)\,dI_\gamma(\chi)\qquad(f\in\Hc(G,K)),
\end{equation}
with $dI_\gamma$ a measure, and if so which measure.

\paragraph{Route A.} The companion paper \cite{JC26} (a revision of \cite{Jo92}) attacks (\ref{e:main}) directly. It truncates the integral over $g$ to a finite union $G(T)$ of double cosets, and expands the spherical function in matrix coefficients. It then uses Casselman's asymptotics and Brion's theorem to write the truncated kernel $I^T(\chi,\gamma)$ as an exponential polynomial in $T$. The $T$-independent part $I(\chi,\gamma)$ is an explicit rational function of $\chi$. Under a regularity hypothesis (R) on the oscillating terms, \cite[Thm.~4.5]{JC26} proves $dI_\gamma=I(\chi,\gamma)\,d\mu$. The hypothesis holds for $PGL(2)$. It fails for $SL(2)$, where $dI_\gamma$ has an atom at the unramified quadratic character. It also fails for $PGL(3)$, where nevertheless $dI_\gamma$ is absolutely continuous \cite[App.~A]{JC26}. So (R) is not the right condition, and \cite[Conj.~4.6]{JC26} predicts the general shape of $dI_\gamma$.

\paragraph{Route B.} Here we derive (\ref{e:main}) from Arthur's local trace formula \cite{Ar89a,Ar91}, in the form given by his theory of elliptic tempered characters \cite{Ar93}. Arthur's formula is an identity of two expansions of one distribution on $G(F)\times G(F)$. For $f=f_1\otimes f_2$ it reads, schematically,
\[
\sum_M c_M\int_{\Gamma_{\ellr}(M)}J_M(\gamma,f)\,d\gamma=\sum_M c_M\int_{\Pi(M)}J_M(\pi,f)\,d\pi,
\]
with weighted orbital integrals on the left and weighted characters on the right. Take $f_1=f$ spherical and $f_2=\phi$ supported on the regular elliptic set. Then almost everything vanishes.
\begin{itemize}
\item On the geometric side only $M=G$ survives, and it is the elliptic orbital integral $I(f,\phi)$ (\S\ref{s:geom}).
\item On the spectral side, a weighted character of the spherical function $f$ vanishes unless its weight is trivial (Lemma~\ref{l:sphvanish}). This rests on the fact that the canonically normalized unramified intertwining operators fix the spherical vector. A character of $\phi$ vanishes on properly induced representations. What remains is a sum over the unitary unramified characters $\chi$ and over the finitely many elliptic tempered virtual characters with $K$-fixed vectors (\S\ref{s:spec}).
\end{itemize}

\paragraph{Results.} The statements below are proved here modulo the normalizations listed in \S\ref{s:status}, which are being checked.

\begin{introtheorem}[existence]\label{t:A}
For every $\phi\in C_c^\infty(G_{\ellr})$ the linear form $f^\vee\mapsto I(f,\phi)$ is given by a finite complex measure $dI_\phi$ on $\Thu/W$.
\end{introtheorem}

\begin{introtheorem}[the absolutely continuous part]\label{t:B}
$dI_\phi=I_B(\chi,\phi)\,d\mu(\chi)+dI^{\mathrm{sing}}_\phi$, where
\[
I_B(\chi,\phi)\,\rho(\chi)=c_G\,J_{M_0}(\chi^{-1},\phi),\qquad c_G=(-1)^{\dim\fa_0^G}|W|^{-1},
\]
$\rho$ is the Plancherel density, and $J_{M_0}$ is Arthur's weighted character for the canonically normalized unramified intertwining operators (\S\ref{s:wchar}). In particular $I_B\rho$ is a smooth function on $\Thu$.
\end{introtheorem}

\begin{introtheorem}[the singular part]\label{t:C}
$dI^{\mathrm{sing}}_\phi$ is a finite sum of smooth measures on subtori of $\Thu/W$, one for each elliptic tempered virtual character, with $K$-fixed vectors, of a Levi subgroup $M\supsetneq M_0$ (\S\ref{s:class}). The subtorus attached to $M$ has dimension $\dim\fa_M^G$. If all $R$-groups of the unitary unramified principal series are trivial, for instance if $Z(G)$ is connected, then $dI^{\mathrm{sing}}_\phi=0$.
\end{introtheorem}

\begin{introtheorem}[$SL(n)$]\label{t:D}
For $G=SL(n)$, $dI^{\mathrm{sing}}_\gamma$ is the sum over the non-trivial unramified $\kappa$ of the push-forwards of the stable spectral measures of the elliptic endoscopic groups, as in \cite[Conj.~4.6(c)]{JC26}.
\end{introtheorem}

The comparison with Route A is Conjecture~\ref{c:AB}: the function $I_B$ of Theorem~\ref{t:B} equals the $T$-independent part $I$ of the truncated kernel of \cite{JC26}. If it holds, Route A gives closed formulas for the density, and Route B explains why it is smooth and what replaces it where (R) fails.

\paragraph{Why this is useful.} Arthur's local trace formula is one of the deepest tools of local harmonic analysis, and most of its terms are hard to compute. In the case treated here every surviving term can be made explicit. The paper is meant both as an independent proof of (\ref{e:main}) and as an introduction to Arthur's local theory through a concrete case.

\paragraph{Plan.} \S\ref{s:lit} reviews the literature before and after 1997. \S\ref{s:bg} collects the unramified background: normalized intertwining operators, $R$-groups, and weighted characters. \S\ref{s:ltf} states Arthur's local trace formula. \S\ref{s:geom} and \S\ref{s:spec} specialize its two sides. \S\ref{s:main} assembles Theorems \ref{t:A}--\ref{t:C}. \S\ref{s:compare} compares the result with Route A. \S\ref{s:endo} treats $SL(n)$. \S\ref{s:ex} works out examples. \S\ref{s:summary} summarizes the results and their status. \S\ref{s:status} lists what remains to be checked, and Appendix~\ref{s:gloss} is a glossary.

% =====================================================================
\section{Literature review}\label{s:lit}

The first version of the companion paper was published in 1992 \cite{Jo92} and corrected in 1997. Arthur's local trace formula and his theory of elliptic tempered characters appeared in between. So we divide the literature at 1997. \S\ref{s:litpre} describes the tools available when the earlier versions were written, and \S\ref{s:litpost} describes what came later and bears on Route~B.

\subsection{Before 1997}\label{s:litpre}

This part describes the tools that were available when the earlier versions of the companion paper were written. They are the spherical theory, intertwining operators and $R$-groups, Arthur's trace formula, and elliptic harmonic analysis.

\paragraph{Spherical harmonic analysis.} Satake identified the spherical Hecke algebra with $W$-invariant functions on the torus of unramified characters \cite{Sa63}. Macdonald gave the explicit formulas for the spherical function and the spherical Plancherel measure \cite{Ma71}. Casselman re-derived Macdonald's formula from intertwining operators and the asymptotics of matrix coefficients \cite{Ca80}, using the theory of Jacquet modules in his notes \cite{Ca95}. Cartier's survey \cite{Ca79} collects the unramified theory. Bernstein and Zelevinsky's geometric lemma \cite{BZ77} describes Jacquet modules of induced representations.

\paragraph{Intertwining operators and $R$-groups.} Knapp and Stein introduced the $R$-group governing the reducibility of induced representations for real groups \cite{KS71}, and Silberger proved the $p$-adic analogue of their dimension theorem \cite{Si78}. Keys computed the $R$-groups of the unramified principal series of Chevalley groups \cite{Ke82}. Langlands' normalizing factors for intertwining operators, built from $L$-functions, go back to his work on Eisenstein series \cite{La76}. Van Dijk's formula for induced characters \cite{vD72} gives the vanishing of characters of cuspidal functions on properly induced representations.

\paragraph{Arthur's trace formula and $(G,M)$-families.} The invariant trace formula \cite{Ar81,Ar88} introduced $(G,M)$-families, weighted orbital integrals and weighted characters, and the splitting formulas. Arthur studied weighted characters as functions of the representation in \cite{Ar89b}. The local trace formula was announced in \cite{Ar89a} and proved in \cite{Ar91}. Arthur then reorganized its spectral side in terms of elliptic tempered virtual characters \cite{Ar93}, and studied the Fourier transforms of weighted orbital integrals \cite{Ar94}. His paper on local character relations \cite{Ar96} relates these objects to endoscopic transfer.

\paragraph{Elliptic harmonic analysis and endoscopy.} Kazhdan's cuspidal geometry \cite{Ka86} and the trace Paley--Wiener theorem of Bernstein, Deligne and Kazhdan \cite{BDK86} control the elliptic part of harmonic analysis. Labesse and Langlands worked out $L$-indistinguishability and $\kappa$-orbital integrals for $SL(2)$ \cite{LL79}, and Langlands and Shelstad defined transfer factors \cite{LS87}. Waldspurger computed the spherical orbital integrals of $GL(N)$ spectrally \cite{Wa89a}, by the truncation method that Route~A follows.

\subsection{After 1997}\label{s:litpost}

This part describes later work that bears on Route~B. It covers normalizations and transfer, the fundamental lemma, elliptic representations, the Plancherel formula, and variants of the local trace formula.

\paragraph{Normalization and transfer of weighted characters.} Arthur proposed a canonical normalization of weighted characters, and a conjectural transfer identity for them \cite{Ar98}. He then studied the transfer of weighted orbital integrals \cite{Ar99}, and stabilized the trace formula in a series of three papers \cite{Ar02,Ar01,Ar03}. The stabilization was completed for twisted groups by M\oe glin and Waldspurger \cite{MW16}. Arthur's classification of representations of classical groups \cite{Ar13} uses the stabilized local and global trace formulas.

\paragraph{The fundamental lemma.} Ng\^o's proof of the fundamental lemma \cite{Ng10}, combined with earlier reductions (see \cite[Literature review]{JC26}), made the endoscopic statements of \S\ref{s:endo} unconditional at the level of orbital integrals.

\paragraph{Elliptic representations.} Reeder related the Euler--Poincar\'e pairing of Iwahori-spherical representations to elliptic representations of Weyl groups \cite{Re01}. Opdam and Solleveld computed Ext groups of tempered representations of affine Hecke algebras, recovering Arthur's elliptic pairing \cite{OS13}. These results describe exactly the elliptic virtual characters with Iwahori-fixed vectors, a class containing those that appear in Theorem~\ref{t:C}.

\paragraph{The Plancherel formula.} Waldspurger's exposition of Harish-Chandra's Plancherel formula \cite{Wa03} supplies complete proofs of the harmonic analysis behind Arthur's spectral side.

\paragraph{Variants of the local trace formula.} Waldspurger proved a local trace formula adapted to the Gross--Prasad conjecture \cite{Wa10}. Beuzart-Plessis proved one for the unitary Gan--Gross--Prasad conjecture \cite{BP20}. Peng made the spectral side of the stable local trace formula explicit when one test function is cuspidal \cite{Pe16}. H.~Li proved a local trace formula for infinitesimal symmetric spaces \cite{Li21}.

\paragraph{Route~A.} The companion paper \cite{JC26} revises \cite{Jo92} by the truncation method, with the regularity hypothesis (R), and with examples in ranks one and two.

\begin{status}
The entries cited in this section have not yet all been checked against publisher records. The page ranges of \cite{Ar01,Ar02,Ar03}, \cite{MW16} and \cite{Wa10} in particular are to be verified.
\end{status}

% =====================================================================
\section{Unramified background}\label{s:bg}

This section collects the facts about unramified principal series that the specialization in \S\S\ref{s:geom}--\ref{s:spec} needs. The essential point is Lemma~\ref{l:sphvanish}: with the canonical unramified normalization, a weighted character of a spherical function has no weight.

\subsection{Notation}\label{s:notation}

This subsection fixes the notation for the group, its parabolic subgroups and the unramified principal series. Example~\ref{ex:notation} makes the notation explicit in rank one, and Table~\ref{tab:notation} at the end of the subsection collects the symbols used in the paper.

We use the notation of \cite[\S2]{JC26}. $A$ is a maximal split torus, $M_0=Z_G(A)$, $B=M_0N_0$ a Borel subgroup, $W$ the Weyl group, $\fa_0=X_*(A)\otimes\mathbb{R}$, and $\cP(M)$ the set of parabolic subgroups with Levi component $M$. For a character $\chi$ of $M_0(F)$ and $P\in\cP(M_0)$, $I_P(\chi)$ is the normalized induced representation, realized on functions on $K$ (the compact picture). $\Phi_K$ is the spherical vector, equal to $1$ on $K$. For unramified $\chi$ write $\chi_\alpha=\chi\circ\alpha^\vee$, an unramified character of $F^\times$, and $L(s,\chi_\alpha)=(1-\chi_\alpha(\pi)q^{-s})^{-1}$. To keep formulas short we assume that $G$ is split. The quasi-split unramified case changes only the local factors attached to roots.

\begin{example}[rank one]\label{ex:notation}
Let $G=SL(2)$ or $PGL(2)$, with $A$ the diagonal torus and $B$ the upper triangular Borel subgroup. Then $M_0=A$, and $\cP(A)=\{B,\bar B\}$ with $\bar B$ the lower triangular Borel subgroup. The Weyl group is $W=\{1,w\}$, the space $\fa_0$ is a line, and there is one reduced positive root $\alpha$.

For $G=SL(2)$ an unramified character of $A(F)$ has the form $\chi(\diag(t,t^{-1}))=\mu(t)$, with $\mu$ an unramified character of $F^\times$. Put $z=\mu(\pi)$, where $\pi$ is a uniformizer of $F$. Here $\alpha^\vee(t)=\diag(t,t^{-1})$, so $\chi_\alpha=\mu$ and $\chi_\alpha(\pi)=z$.

For $G=PGL(2)$ an unramified character of $A(F)$ has the form $\chi(\diag(a,b))=\mu(a/b)$, and again we put $z=\mu(\pi)$. Here $\alpha^\vee(t)$ is the image of $\diag(t,t^{-1})$, so $\chi_\alpha=\mu^2$ and $\chi_\alpha(\pi)=z^2$.

In both cases $\chi$ is unitary exactly when $|z|=1$, so $\Thu$ is the unit circle. The element $w$ sends $z$ to $z^{-1}$, and a Satake transform $f^\vee$ is a Laurent polynomial in $z$ that is invariant under $z\mapsto z^{-1}$. These are the coordinates of \cite[\S6.0]{JC26}.
\end{example}

{\small
\begin{longtable}{@{}>{\raggedright\arraybackslash}p{0.27\textwidth}>{\raggedright\arraybackslash}p{0.69\textwidth}@{}}
\caption{Notation. Symbols taken from the companion paper are defined there, at the place cited.}\label{tab:notation}\\
\toprule
Symbol & Meaning\\
\midrule
\endfirsthead
\toprule
Symbol & Meaning\\
\midrule
\endhead
\bottomrule
\endfoot
$F$, $\OO$, $q$ & A $p$-adic field, its ring of integers, and the order of its residue field (\S\ref{s:intro}).\\
$\pi$ & A uniformizer of $F$ when it is the argument of a character, as in $\chi_\alpha(\pi)$. Elsewhere $\pi$, $\sigma$ and $\tau$ denote representations or virtual characters.\\
$G$, $K$ & A connected unramified semisimple group over $F$, and a hyperspecial maximal compact subgroup of $G(F)$.\\
$A$, $M_0$, $B$, $W$ & A maximal split torus, its centralizer, a Borel subgroup $B=M_0N_0$, and the Weyl group.\\
$\cL$, $\cL(M)$, $\cP(M)$, $\cP^L(M)$ & The Levi subgroups containing $M_0$. Those containing $M$. The parabolic subgroups with Levi component $M$. The same inside a Levi subgroup $L$.\\
$\fa_M$, $\fa_M^L$, $\fa_0$ & $\fa_M=X_*(A_M)\otimes\mathbb{R}$, with $A_M$ the split part of the center of $M$. $\fa_M^L$ is the kernel of $\fa_M\to\fa_L$. $\fa_0=\fa_{M_0}$. Since $G$ is semisimple, $\fa_0^G=\fa_0$.\\
$H_M$, $H_P$ & The map $M(F)\to\fa_M$ with $e^{\ip{H_M(m),\lambda}}=|\lambda(m)|$ for rational characters $\lambda$ of $M$, and its extension to $G(F)$ by $H_P(mnk)=H_M(m)$ \cite[Glossary]{JC26}.\\
$\Thu$, $\chi_\Lambda$, $\chi_\alpha$ & The compact torus of unitary unramified characters of $M_0(F)$. The twist $\chi\cdot e^{\Lambda(H_{M_0}(\cdot))}$. The character $\chi\circ\alpha^\vee$ of $F^\times$.\\
$\Hc(G,K)$, $f^\vee$ & The spherical Hecke algebra and the Satake transform.\\
$d\mu$, $\rho$, $d\chi$ & The Plancherel measure on $\Thu/W$, its density, and Haar measure on $\Thu$.\\
$I_P(\chi)$, $\Phi_K$ & The normalized induced representation in the compact picture, and its spherical vector.\\
$J_{P'|P}$, $r_{P'|P}$, $R_{P'|P}$, $c_\alpha$ & The standard intertwining operator, the normalizing factor, the normalized operator, and the rank-one $c$-function (\S\ref{s:norm}).\\
$\theta_Q(\Lambda)$ & The product of the $\Lambda(\alpha^\vee)$ over the simple roots $\alpha$ of $Q$, divided by a covolume \cite[\S6]{Ar81}, \cite[(2.16)]{JC26}.\\
$\mathcal R_M(\chi,P)$, $J_M(\chi,f)$, $J^L_M$ & The limit of the $(G,M)$-family of normalized operators, the weighted character, and its analogue for a Levi subgroup $L\supseteq M$ (\S\ref{s:wchar}).\\
$G_{\ellr}$, $O_\gamma$, $I(f,\phi)$ & The regular elliptic set, the orbital integral, and the distribution of Proposition~\ref{p:geom}.\\
$D(\gamma)$ & The Weyl discriminant $\det(1-\Ad(\gamma))$ on $\mathfrak g/\mathfrak g_\gamma$.\\
$K_f$, $u(x,T)$, $J^T(f)$, $J(f)$ & The kernel, Arthur's truncation function, the truncated integral and its constant term (\S\ref{s:trunc}).\\
$J_M(\gamma,f)$, $v_M$, $\Gamma_{\ellr}(M)$ & The weighted orbital integral, its weight, and the elliptic regular classes of $M(F)$ (\S\ref{s:expansions}).\\
$\Pi_{\mathrm{temp}}(M)$, $\Tdisc(M)$, $\Tell(M)$ & The tempered representations of $M(F)$, and Arthur's sets of virtual characters (\S\ref{s:ellform}).\\
$d_M^G(L_1,L_2)$ & The constants of the splitting formula (\ref{e:split}) \cite{Ar88}.\\
$W_\chi$, $W^\circ_\chi$, $R_\chi$, $R^M_\chi$ & The stabilizer of $\chi$ in $W$, its reflection subgroup, the $R$-group, and the $R$-group relative to $M$ (\S\ref{s:R}).\\
$\chi_0$ & The unramified quadratic character of $SL(2)$, the point $z=-1$ of Example~\ref{ex:notation}.\\
$dI_\phi$, $dI_\gamma$, $dI^{\mathrm{sing}}_\phi$, $I_B$, $c_G$ & The spectral measure of $\phi$ or of $\gamma$, its singular part, the density of its absolutely continuous part, and the constant of Theorem~\ref{t:B}.\\
$c_{M,\tau}$, $c_\tau$ & Constants from \cite{Ar93} in Theorem~\ref{t:main}. They are not determined in this paper.\\
$L_M^\vee$ & A lattice in $i\fa_M^*$ such that the twist $\tau_\Lambda$ depends only on $\Lambda$ modulo $L_M^\vee$ (Theorem~\ref{t:main}).\\
$G(T)$, $I^T$, $I(\chi,\gamma)$, (R), (R$'$) & Route~A: the truncation set, the truncated kernel, its $T$-independent part, and the regularity hypotheses \cite[\S\S3--4, Rem.~2026.20]{JC26}.\\
$\delta_B$, $Q_a$, $\Gamma_\chi$ & The modulus character of $B$, the constant in the volume formula for $KaK$ \cite[\S3.2]{JC26}, and the spherical function.\\
$Q$, (U$_r$), $B_r$, $K_1$ & Rank-one notation of \cite[\S6]{JC26}: $Q=1+q^{-1}$, the type of an unramified elliptic element of depth $r$, the ball of vertices that it fixes, and the first principal congruence subgroup.\\
$\kappa$, $H_m$, $\Delta_\kappa$, $c_S$ & For $SL(n)$: an unramified character of $PGL(n,F)$, the elliptic endoscopic group, the transfer factor and a constant \cite[\S5]{JC26}.\\
\end{longtable}
}

\subsection{Normalized intertwining operators}\label{s:norm}

This subsection recalls the standard intertwining operators between unramified principal series and fixes their normalization. The point is (\ref{e:fix}): the normalized operators fix the spherical vector. Everything in this paper that is special to spherical functions comes from this fact.

For $P,P'\in\cP(M_0)$ let $J_{P'|P}(\chi):I_P(\chi)\to I_{P'}(\chi)$ be the standard intertwining operator, defined by a convergent integral for $\chi$ in a cone and by meromorphic continuation in general. Gindikin--Karpelevich and Casselman \cite[Thm.~3.1]{Ca80} give, in the compact picture,
\begin{equation}\label{e:GK}
J_{P'|P}(\chi)\Phi_K=\prod_{\alpha}c_\alpha(\chi)\,\Phi_K,\qquad c_\alpha(\chi)=\frac{1-q^{-1}\chi_\alpha(\pi)}{1-\chi_\alpha(\pi)}=\frac{L(0,\chi_\alpha)}{L(1,\chi_\alpha)},
\end{equation}
the product over the reduced roots $\alpha$ that are positive for $P$ and negative for $P'$. Define
\begin{equation}\label{e:rnorm}
r_{P'|P}(\chi)=\prod_\alpha\frac{L(0,\chi_\alpha)}{L(1,\chi_\alpha)},\qquad R_{P'|P}(\chi)=r_{P'|P}(\chi)^{-1}J_{P'|P}(\chi).
\end{equation}
For unramified data these are Langlands' normalizing factors \cite{La76}, with the $\varepsilon$-factors equal to $1$ for an additive character of conductor $\OO$. They satisfy Arthur's conditions on normalizing factors: multiplicativity, $R_{P|P'}R_{P'|P}=1$, unitarity and holomorphy on the unitary axis, and compatibility with the $K$-action. By (\ref{e:GK}),
\begin{equation}\label{e:fix}
R_{P'|P}(\chi)\Phi_K=\Phi_K\qquad\text{for all }P,P'\in\cP(M_0)\text{ and all }\chi.
\end{equation}

\begin{status}
For split $G$ this is standard. Arthur's conditions on normalizing factors, and his canonical normalization \cite{Ar98}, have not yet been compared in detail with (\ref{e:rnorm}); see \S\ref{s:status}, item~\ref{n:1}.
\end{status}

\begin{example}[rank one]\label{ex:norm}
In the setting of Example~\ref{ex:notation}, put $y=\chi_\alpha(\pi)$, so that $y=z$ for $SL(2)$ and $y=z^2$ for $PGL(2)$. There is one pair of distinct parabolic subgroups, and (\ref{e:GK}) reads
\[
J_{\bar B|B}(\chi)\Phi_K=c_\alpha(\chi)\Phi_K,\qquad c_\alpha(\chi)=\frac{1-q^{-1}y}{1-y}.
\]
On the unit circle $c_\alpha$ has no zero, and it has a pole exactly at $y=1$. This is the point $z=1$ for $SL(2)$, and the two points $z=\pm1$ for $PGL(2)$. By (\ref{e:rnorm}) the normalizing factor is $r_{\bar B|B}(\chi)=c_\alpha(\chi)$. So the normalized operator $R_{\bar B|B}(\chi)=c_\alpha(\chi)^{-1}J_{\bar B|B}(\chi)$ satisfies $R_{\bar B|B}(\chi)\Phi_K=\Phi_K$ for all $\chi$, which is (\ref{e:fix}). The pole of the unnormalized operator on the spherical vector has been removed.
\end{example}

\subsection{Weighted characters}\label{s:wchar}

This subsection defines Arthur's weighted characters. They are the terms of the spectral side of his trace formula, and they replace ordinary characters when a Levi subgroup $M\ne G$ is involved. The subsection then proves Lemma~\ref{l:sphvanish}, on which the rest of the paper depends: a weighted character of a spherical function vanishes unless its weight is trivial.

Fix $P\in\cP(M_0)$ and unitary $\chi$. For $Q\in\cP(M_0)$ and $\Lambda\in i\fa_0^*$ put
\[
\mathcal R_Q(\Lambda,\chi,P)=R_{Q|P}(\chi)^{-1}R_{Q|P}(\chi_\Lambda),
\]
where $\chi_\Lambda=\chi\cdot e^{\Lambda(H_{M_0}(\cdot))}$. Arthur shows that $\{\mathcal R_Q(\Lambda,\chi,P)\}_Q$ is a $(G,M_0)$-family of operators \cite[\S7]{Ar81}. More generally, for $M\supseteq M_0$ and a Levi $L\supseteq M$ there is an operator $\mathcal R^L_M(\chi,P)$, the limit at $\Lambda=0$ of the associated sum over $\cP^L(M)$ \cite[\S6]{Ar81}. The \emph{weighted character} of $f\in C_c^\infty(G(F))$ is
\[
J_M(\chi,f)=\tr\big(\mathcal R_M(\chi,P)\,I_P(\chi,f)\big).
\]
It does not depend on $P$. For $M=G$ it is the character $\tr I_P(\chi,f)$. The same construction applies to every tempered representation of a Levi subgroup.

\begin{lemma}\label{l:sphvanish}
Let $f\in\Hc(G,K)$, let $\chi$ be unramified, and let $M\subsetneq L$ be Levi subgroups containing $M_0$. Then $J^L_M(\chi,f)=0$, where $J^L_M$ is formed with the normalization (\ref{e:rnorm}).
\end{lemma}

\begin{proof}
Since $f$ is bi-$K$-invariant, $I_P(\chi,f)$ maps $I_P(\chi)$ into the line $\mathbb{C}\Phi_K$ and kills the $K$-invariant complement of that line. So $I_P(\chi,f)=f^\vee(\chi)\,e_K$, with $e_K$ the projection onto $\mathbb{C}\Phi_K$, and $J^L_M(\chi,f)=f^\vee(\chi)\ip{\mathcal R^L_M(\chi,P)\Phi_K,\Phi_K^\vee}$. By (\ref{e:fix}), $\mathcal R_Q(\Lambda,\chi,P)\Phi_K=\Phi_K$ for every $Q$ and $\Lambda$. So on $\Phi_K$ the family is the constant $(L,M)$-family $1$. The limit attached to a constant $(L,M)$-family is $0$ when $M\ne L$, because $\sum_{Q\in\cP^L(M)}\theta_Q(\Lambda)^{-1}=0$ \cite[Lemma~6.2]{Ar81}. Hence $\mathcal R^L_M(\chi,P)\Phi_K=0$.
\end{proof}

\begin{remark}
Lemma~\ref{l:sphvanish} depends on the normalization. With unnormalized operators, the weighted characters of spherical functions would carry the logarithmic derivatives of the $c$-functions. These are exactly the terms that appear in the $T$-independent part of the truncated kernel of Route~A; compare \S\ref{s:compare}.
\end{remark}

\begin{example}[a weighted character in rank one]\label{ex:wchar}
In the setting of Example~\ref{ex:notation}, take $M=M_0=A$ and $P=B$. The $(G,A)$-family has two members,
\[
\mathcal R_B(\Lambda,\chi,B)=1,\qquad \mathcal R_{\bar B}(\Lambda,\chi,B)=R_{\bar B|B}(\chi)^{-1}R_{\bar B|B}(\chi_\Lambda).
\]
They agree at $\Lambda=0$, as the members of a $(G,A)$-family must. Write $\xi=\Lambda(\alpha^\vee)$. With the conventions of \cite[Example~2.5]{JC26}, $\theta_B(\Lambda)=v^{-1}\xi$ and $\theta_{\bar B}(\Lambda)=-v^{-1}\xi$, where $v$ is the volume of $\fa_0/\mathbb{Z}\alpha^\vee$. So
\[
\mathcal R_A(\chi,B)=\lim_{\xi\to0}\,v\,\frac{1-R_{\bar B|B}(\chi)^{-1}R_{\bar B|B}(\chi_\Lambda)}{\xi}=-v\,R_{\bar B|B}(\chi)^{-1}\,\frac{d}{d\xi}R_{\bar B|B}(\chi_\Lambda)\Big|_{\xi=0}.
\]
The operator $\mathcal R_A(\chi,B)$ is thus a logarithmic derivative of the normalized intertwining operator, and $J_A(\chi,f)=\tr(\mathcal R_A(\chi,B)I_B(\chi,f))$. By contrast, $J_G(\chi,f)=\tr I_B(\chi,f)$ is the ordinary character.

On the spherical vector, Example~\ref{ex:norm} gives $R_{\bar B|B}(\chi_\Lambda)\Phi_K=\Phi_K$ for all $\Lambda$. The difference quotient is then zero, so $\mathcal R_A(\chi,B)\Phi_K=0$. This is Lemma~\ref{l:sphvanish} for $M=A$ and $L=G$.
\end{example}

\subsection{\texorpdfstring{$R$}{R}-groups and the unramified elliptic set}\label{s:R}

This subsection recalls the Knapp--Stein $R$-group of an unramified principal series, and the notion of an elliptic element of it. The $R$-group controls where $I_P(\chi)$ is reducible. It is needed to describe the singular part of the spectral measure in Theorem~\ref{t:C}.

For unitary unramified $\chi$ let $W_\chi=\{w\in W:w\chi=\chi\}$, and let $W^\circ_\chi$ be the subgroup generated by the reflections $s_\alpha$ with $\chi_\alpha=1$. The Knapp--Stein $R$-group is $R_\chi=W_\chi/W^\circ_\chi$. It governs the decomposition of $I_P(\chi)$: the commuting algebra is a twisted group algebra of $R_\chi$, which for unramified principal series is untwisted \cite{Ke82}. For Chevalley groups Keys computed $R_\chi$ \cite{Ke82}. In particular:
\begin{itemize}
\item if the derived group of the dual group $\hat G$ is simply connected (equivalently $Z(G)$ is connected), then $R_\chi=1$ for all unitary unramified $\chi$;
\item for $G=SL(n)$, $R_\chi$ is cyclic of order dividing $n$, and $R_\chi\ne1$ exactly on a finite union of subtori of $\Thu$. For $SL(2)$ this is the point $\chi_0$, the unramified quadratic character.
\end{itemize}

An element $r\in R_\chi$, represented by $w\in W_\chi$, is \emph{elliptic} if $w$ has no non-zero fixed vector in $\fa_0^G$. Arthur's elliptic tempered virtual characters with $K$-fixed vectors are exactly the virtual characters $\tau=(\chi,r)$ with $r$ elliptic (\S\ref{s:class}).

\begin{example}[$R$-groups in rank one]\label{ex:R}
In the setting of Example~\ref{ex:notation}, $w\chi=\chi^{-1}$. So $W_\chi=W$ exactly when $z=\pm1$, and $W_\chi=1$ otherwise.

Let $G=SL(2)$. For $z=1$ we have $\chi_\alpha=1$, so $W^\circ_\chi=W$ and $R_\chi=1$. For $z=-1$, which is the character $\chi_0$, we have $\chi_\alpha=\mu\ne1$, so $W^\circ_{\chi_0}=1$ and $R_{\chi_0}=W$ has order $2$. Its non-trivial element acts by $-1$ on the line $\fa_0$ and has no non-zero fixed vector, so it is elliptic. This is the only unitary unramified character of $SL(2)$ with a non-trivial $R$-group, and \S\ref{s:sl2} shows what it contributes.

Let $G=PGL(2)$. For $z=\pm1$ we have $\chi_\alpha(\pi)=z^2=1$, so $\chi_\alpha=1$, $W^\circ_\chi=W$ and $R_\chi=1$. All $R$-groups are trivial. This agrees with the first item of the list above, since the center of $PGL(2)$ is trivial.
\end{example}

% =====================================================================
\section{Arthur's local trace formula}\label{s:ltf}

This section states the parts of Arthur's local trace formula that we use. We quote; nothing here is new.

\subsection{The truncated kernel}\label{s:trunc}

This subsection recalls how Arthur defines the distribution $J(f)$. The kernel of a convolution operator is not integrable on the diagonal, so it is truncated first.

The group $G(F)\times G(F)$ acts on $L^2(G(F))$ by $(R(x_1,x_2)\psi)(y)=\psi(x_1^{-1}yx_2)$. For $f=f_1\otimes f_2$ the operator $R(f)$ has kernel
\[
K_f(x,y)=\int_{G(F)}f_1(u)f_2(x^{-1}uy)\,du.
\]
Its restriction to the diagonal is not integrable in general. Arthur multiplies it by the characteristic function $u(x,T)$ of a large compact subset of $G(F)$, defined by a point $T$ in the positive chamber through the convex hull of $W\cdot T$. He shows that $J^T(f)=\int_{G(F)}K_f(x,x)u(x,T)\,dx$ is asymptotic to a polynomial in $T$ \cite{Ar89a,Ar91}. Its constant term $J(f)$ is a distribution on $G(F)\times G(F)$.

\subsection{The geometric and spectral expansions}\label{s:expansions}

This subsection states the two expansions of $J(f)$. Their equality is the local trace formula. The geometric expansion is a sum of weighted orbital integrals, and the spectral expansion is a sum of weighted characters.

Arthur proves two expansions of $J(f)$ \cite[Props.~6.1 and~11.3, Thm.~12.2]{Ar91}:
\begin{align}
J(f)&=\sum_{M\in\cL}\frac{|W_0^M|}{|W_0^G|}(-1)^{\dim\fa_M^G}\int_{\Gamma_{\ellr}(M)}J_M(\gamma,f)\,d\gamma,\label{e:geo}\\
J(f)&=\sum_{M\in\cL}\frac{|W_0^M|}{|W_0^G|}(-1)^{\dim\fa_M^G}\int_{\Pi_{\mathrm{temp}}(M)}J_M(\pi,f)\,d\pi.\label{e:spec}
\end{align}
Here $\cL$ is the set of Levi subgroups containing $M_0$, and $\Gamma_{\ellr}(M)$ is the set of elliptic regular conjugacy classes of $M(F)$. The weighted orbital integral $J_M(\gamma,f)$ is
\[
J_M(\gamma,f)=|D(\gamma)|\int_{M_\gamma\backslash G}\int_{M_\gamma\backslash G}f_1(x_1^{-1}\gamma x_1)f_2(x_2^{-1}\gamma x_2)\,v_M(x_1,x_2)\,dx_1\,dx_2,
\]
with $v_M$ the weight of the $(G,M)$-family $e^{-\Lambda(H_P(x_1)-H_{\bar P}(x_2))}$. In (\ref{e:spec}), $J_M(\pi,f)$ is a weighted character of $G(F)\times G(F)$ attached to $\pi\otimes\pi^\vee$. The measure $d\pi$ and the precise domain of integration are as in \cite{Ar91}.

\begin{example}[rank one]\label{ex:ltf}
Let $G=SL(2)$ or $PGL(2)$. Then $\cL=\{A,G\}$. For $M=G$ the coefficient in (\ref{e:geo}) and (\ref{e:spec}) is $1$. For $M=A$ we have $|W_0^A|=1$, $|W_0^G|=2$ and $\dim\fa_A^G=1$, so the coefficient is $-\frac12$. The geometric expansion (\ref{e:geo}) therefore reads
\[
J(f)=\int_{\Gamma_{\ellr}(G)}J_G(\gamma,f)\,d\gamma-\frac12\int_{\Gamma_{\ellr}(A)}J_A(\gamma,f)\,d\gamma.
\]
The first term contains ordinary orbital integrals over the elliptic classes of $G(F)$. The second contains weighted orbital integrals over elements of the split torus $A(F)$, since every element of a torus is elliptic in that torus. The spectral expansion (\ref{e:spec}) has the same two coefficients. The coefficient $-\frac12$ is the constant $c_G$ of Theorem~\ref{t:B} in rank one.
\end{example}

\subsection{The elliptic form}\label{s:ellform}

This subsection recalls Arthur's reorganization of the spectral side by virtual characters. It is the form in which the singular part of the spectral measure will appear.

In \cite{Ar93} Arthur reorganizes the spectral side using the virtual characters $\tau\in\Tdisc(M)$: triples $(L,\sigma,r)$ with $\sigma$ a square-integrable representation of a Levi $L\subseteq M$ and $r$ in the $R$-group of $\sigma$ in $M$, up to equivalence. The elliptic ones, $\Tell(M)$, are those for which $r$ is elliptic in the sense of \S\ref{s:R}. The spectral side becomes a sum over $M$ of integrals over $\Tdisc(M)$, with explicit constants. We need its form for $f=f_1\otimes f_2$ with $f_2$ cuspidal (\S\ref{s:spec}).

\subsection{Splitting}\label{s:splitting}

Our test functions are products $f_1\otimes f_2$. The splitting formula of this subsection expresses a weighted character of a product through weighted characters of the two factors. In \S\ref{s:spec} it lets us remove the spherical factor.

For $f=f_1\otimes f_2$ the weighted characters split \cite[Cor.~7.4]{Ar88}, \cite[\S10]{Ar91}:
\begin{equation}\label{e:split}
J_M(\pi,f_1\otimes f_2)=\sum_{L_1,L_2\in\cL(M)}d_M^G(L_1,L_2)\,J^{L_1}_M(\pi,f_{1,Q_1})\,J^{L_2}_M(\pi^\vee,f_{2,Q_2}).
\end{equation}
Here $d_M^G(L_1,L_2)$ is non-zero only if $\fa_M^{L_1}\oplus\fa_M^{L_2}=\fa_M^G$, $d_M^G(M,G)=1$, $f_{i,Q_i}$ are constant terms, and the $Q_i\in\cP(L_i)$ are fixed by a choice of a generic point. The same holds for weighted orbital integrals.

\begin{status}
The constants in (\ref{e:geo})--(\ref{e:spec}), the exact domain and measure on the spectral side, and the form of (\ref{e:split}) in the local trace formula (as opposed to the invariant trace formula of \cite{Ar88}) are being checked against \cite{Ar91,Ar93}; see \S\ref{s:status}, item~\ref{n:2}.
\end{status}

% =====================================================================
\section{The geometric side}\label{s:geom}

This section specializes the geometric expansion (\ref{e:geo}) to the product of a spherical function and a function supported on the elliptic set. Proposition~\ref{p:geom} shows that only the term $M=G$ survives, and that it is the distribution $I(f,\phi)$ studied in the companion paper.

Let $f_1=f\in\Hc(G,K)$, and let $f_2=\phi\in C_c^\infty(G(F))$ be supported in the regular elliptic set $G_{\ellr}$.

\begin{proposition}\label{p:geom}
With $\mathbf f=f\otimes\phi$:
\begin{enumerate}[label=(\alph*)]
\item $x\mapsto K_{\mathbf f}(x,x)$ has compact support, so $J^T(\mathbf f)$ is constant for $T$ large, and
\[
J(\mathbf f)=\int_{G(F)}\int_{G(F)}f(u)\phi(x^{-1}ux)\,du\,dx=I(f,\phi),
\]
the distribution of \cite[(1.1), (4.8)]{JC26};
\item in (\ref{e:geo}) the terms with $M\ne G$ vanish, and the term $M=G$ is the Weyl integration formula for $I(f,\phi)$.
\end{enumerate}
\end{proposition}

\begin{proof}
(a) For $\phi$ supported on $G_{\ellr}$, the set of $x$ with $\phi(x^{-1}ux)\ne0$ for some $u$ in the support of $f$ is compact, because centralizers of elliptic elements are compact (Kazhdan \cite[\S4]{Ka86}; \cite[Lemma~4.4]{JC26}). So $K_{\mathbf f}(x,x)$ has compact support, and $u(x,T)=1$ on it for $T$ large.

(b) Let $M\ne G$ and $\gamma\in\Gamma_{\ellr}(M)$. If $\gamma$ is regular in $G$, its centralizer in $G$ is a maximal torus of $M$, which contains the non-trivial split torus $A_M$; so $\gamma$ is not elliptic in $G$. Hence $\phi(x^{-1}\gamma x)=0$ for all $x$, and $J_M(\gamma,\mathbf f)=0$. The $\gamma$ that are singular in $G$ form a set of measure zero in $\Gamma_{\ellr}(M)$. For $M=G$ the weight is $1$, and $J_G(\gamma,\mathbf f)=|D(\gamma)|O_\gamma(f)O_\gamma(\phi)$ with the measures of \cite[(4.10)]{JC26}.
\end{proof}

\begin{example}[rank one]\label{ex:geom}
Let $G=SL(2)$ or $PGL(2)$, and consider the term $M=A$ of Example~\ref{ex:ltf}. A regular element $\gamma$ of $A(F)$ has centralizer $A$, which is split. So $\gamma$ is not elliptic in $G$, and neither is any conjugate of it. Since $\phi$ is supported on $G_{\ellr}$, the factor $\phi(x_2^{-1}\gamma x_2)$ in $J_A(\gamma,f\otimes\phi)$ is zero for all $x_2$. The term with coefficient $-\frac12$ therefore vanishes. What remains is the term $M=G$, an integral of $|D(\gamma)|O_\gamma(f)O_\gamma(\phi)$ over the elliptic classes of $G(F)$, and it equals $I(f,\phi)$. This is Proposition~\ref{p:geom} in rank one.
\end{example}

% =====================================================================
\section{The spectral side}\label{s:spec}

This section specializes the spectral expansion to the same test function $f\otimes\phi$. There are three steps. Splitting removes the weight from the spherical factor (\S\ref{s:specsplit}). Cuspidality of $\phi$ removes the properly induced data (\S\ref{s:cusp}). A classification then lists the virtual characters that survive (\S\ref{s:class}).

\subsection{Splitting with a spherical and an elliptic function}\label{s:specsplit}

This subsection applies the splitting formula (\ref{e:split}) and Lemma~\ref{l:sphvanish}. The result is that the spherical function contributes only its Satake transform.

\begin{proposition}\label{p:split}
Let $f\in\Hc(G,K)$ and $\phi\in C_c^\infty(G_{\ellr})$, and let $\pi$ be a tempered representation of $M(F)$ induced from a unitary unramified character $\chi$ of $M_0(F)$, or more generally any $\tau\in\Tdisc(M)$. Then
\[
J_M(\pi,f\otimes\phi)=f^\vee(\chi)\,J_M(\pi^\vee,\phi).
\]
If $\pi$ has no $(K\cap M)$-fixed vector, both sides vanish.
\end{proposition}

\begin{proof}
In (\ref{e:split}), the factor $J^{L_1}_M(\pi,f_{1,Q_1})$ vanishes for $L_1\ne M$ by Lemma~\ref{l:sphvanish}, applied to the constant term $f_{1,Q_1}\in\Hc(L_1,K\cap L_1)$. For $L_1=M$ the condition $\fa_M^M\oplus\fa_M^{L_2}=\fa_M^G$ forces $L_2=G$, and $d^G_M(M,G)=1$. The remaining factor is $J^M_M(\pi,f_{1,Q})=\tr\pi(f_{1,Q})=\tr I_Q(\pi,f)$, which equals $f^\vee(\chi)$ if $\pi$ is a constituent of an unramified principal series with a $K\cap M$-fixed vector, and $0$ if it has none.
\end{proof}

\subsection{Cuspidality of \texorpdfstring{$\phi$}{phi}}\label{s:cusp}

This subsection shows that a function supported on the elliptic set is cuspidal: its trace vanishes on properly induced representations.

\begin{lemma}\label{l:cusp}
If $\phi$ is supported on $G_{\ellr}$, then $\tr\sigma(\phi)=0$ for every representation $\sigma$ of finite length induced from a proper parabolic subgroup.
\end{lemma}

\begin{proof}
By van Dijk's formula \cite{vD72}, the character of a properly induced representation is supported on the conjugates of the Levi subgroup, which meet $G_{\ellr}$ in a set of measure zero.
\end{proof}

So $\phi$ is cuspidal, and in the elliptic form of the spectral side only elliptic data survive in the $\phi$-factor.

\subsection{Which virtual characters survive}\label{s:class}

Combining Proposition~\ref{p:split} and Lemma~\ref{l:cusp} with Arthur's elliptic form, the spectral side of $J(f\otimes\phi)$ is a sum over $M\in\cL$ of integrals, over the virtual characters $\tau\in\Tdisc(M)$ with $K\cap M$-fixed vectors, of $f^\vee(\chi_\tau)J_M(\tau^\vee,\phi)$, where $\chi_\tau$ is the unramified character underlying $\tau$. We now list these $\tau$.

\begin{lemma}\label{l:noDS}
A square-integrable representation of $L(F)$, $L$ a Levi subgroup of $G$ with $\dim\fa_{M_0}^L>0$, has no vector fixed by the hyperspecial subgroup $K\cap L(F)$.
\end{lemma}

\begin{proof}[Sketch]
A $K\cap L$-spherical irreducible representation is a constituent of an unramified principal series $I^L(\chi)$, and its spherical matrix coefficient is the spherical function $\Gamma^L_\chi$. By Macdonald's formula, $|\Gamma^L_\chi(a)|$ is bounded below by a non-zero multiple of $\delta_B^{1/2}(a)$ on a cone in $A^-$, for every $\chi$. By the volume formula $\vol(KaK)=\delta_B(a)^{-1}Q_a$, this is not square-integrable.
\end{proof}

\begin{status}
The lower bound in the sketch needs care at non-regular $\chi$, where the terms of Macdonald's formula can cancel. The statement is presumably in the literature (Casselman's criterion for square-integrability \cite{Ca95} gives it directly, since the exponents of a spherical representation lie in $\delta_B^{1/2}\cdot\Thu$ up to $W$). To be completed.
\end{status}

By Lemma~\ref{l:noDS}, in a triple $\tau=(L,\sigma,r)$ with a $K$-fixed vector we must have $L=M_0$ and $\sigma=\chi$ a unitary unramified character. So:
\begin{itemize}
\item $M=M_0$: $\Tdisc(M_0)=\Thu$, and the contribution is
\[
\frac{(-1)^{\dim\fa_0^G}}{|W|}\int_{\Thu}f^\vee(\chi)\,J_{M_0}(\chi^{-1},\phi)\,d\chi.
\]
\item $M\supsetneq M_0$: $\tau=(M_0,\chi,r)$ with $r\in R^M_\chi$ elliptic in $M$, i.e.\ represented by $w\in W^M_\chi$ without non-zero fixed vector in $\fa_0^M$. Unramified twists by $i\fa_M^*$ act on such $\tau$, so each orbit is a torus of dimension $\dim\fa_M^G$ in $\Thu/W$. The contribution is an integral of $f^\vee(\chi_\tau)J_M(\tau^\vee,\phi)$ over these tori, with the constant of \cite{Ar93}.
\item $M=G$: $\tau=(M_0,\chi,r)$ with $r\in R_\chi$ elliptic in $G$. These $\tau$ form a finite set, and contribute point masses $f^\vee(\chi_\tau)\tr\tau^\vee(\phi)$, with constants.
\end{itemize}

\begin{example}[rank one]\label{ex:spec}
Let $G=PGL(2)$. All $R$-groups are trivial (Example~\ref{ex:R}), so only the first item of the list occurs. The spectral side of $J(f\otimes\phi)$ is
\[
-\frac12\int_{\Thu}f^\vee(\chi)\,J_A(\chi^{-1},\phi)\,d\chi.
\]
Let $G=SL(2)$. The same term occurs, and the third item contributes as well. By Example~\ref{ex:R} the only $\tau=(A,\chi,r)$ with $r$ elliptic is the one at $\chi=\chi_0$, and it gives a point mass at $\chi_0$ (\S\ref{s:sl2}). In rank one there is no Levi subgroup strictly between $A$ and $G$, so no tori of positive dimension occur.
\end{example}

% =====================================================================
\section{The main theorem}\label{s:main}

This section equates the geometric side of \S\ref{s:geom} with the spectral side of \S\ref{s:spec}. The result is Theorem~\ref{t:main}. Theorems~\ref{t:A}, \ref{t:B} and \ref{t:C} of the Introduction follow from it.

\begin{theorem}\label{t:main}
Let $\phi\in C_c^\infty(G_{\ellr})$. With the normalizations of \S\ref{s:status} assumed, for all $f\in\Hc(G,K)$,
\[
I(f,\phi)=c_G\int_{\Thu}f^\vee(\chi)\,J_{M_0}(\chi^{-1},\phi)\,d\chi+\sum_{M\supsetneq M_0}\sum_{[\tau]}c_{M,\tau}\int_{i\fa_M^*/L_M^\vee}f^\vee(\chi_{\tau,\Lambda})\,J_M(\tau_\Lambda^\vee,\phi)\,d\Lambda,
\]
where $[\tau]$ runs over the finitely many twist-orbits of $\tau=(M_0,\chi,r)$ with $r$ elliptic in $M$, and the $c_{M,\tau}$ are constants from \cite{Ar93}. Theorems~\ref{t:A}, \ref{t:B} and \ref{t:C} follow.
\end{theorem}

\begin{proof}[Proof outline]
Equate Proposition~\ref{p:geom} with the spectral side of \S\ref{s:class}.
\begin{itemize}
\item Theorem~\ref{t:B}: the normalized intertwining operators are holomorphic and unitary on the unitary axis, so $\chi\mapsto\mathcal R_{M_0}(\chi,P)$ is smooth on $\Thu$. Since $I_P(\chi,\phi)$ has finite rank and depends smoothly on $\chi$ in the compact picture, $J_{M_0}(\chi^{-1},\phi)$ is smooth. Divide by $\rho$ to get $I_B$.
\item Theorem~\ref{t:C}: the same argument on each torus gives smooth densities. The list of \S\ref{s:class} gives the supports.
\item Theorem~\ref{t:A}: there are finitely many components, each with a smooth density on a compact torus, so the total variation is finite.
\end{itemize}
\end{proof}

\begin{remark}
Theorem~\ref{t:B} shows that $I_B\rho$ is always smooth. In Route~A, the $T$-independent part $I$ has double poles on the root walls, canceled by the double zeros of $\rho$ \cite[Rem.~6.2, (A.4)]{JC26}. Route B explains this cancellation: it is the smoothness of a weighted character.
\end{remark}

\begin{example}[rank one]\label{ex:main}
Let $G=PGL(2)$. With the same normalizations assumed, Theorem~\ref{t:main} and Example~\ref{ex:spec} give
\[
I(f,\phi)=-\frac12\int_{\Thu}f^\vee(\chi)\,J_A(\chi^{-1},\phi)\,d\chi.
\]
There is no singular part, and $I_B\rho=-\frac12J_A(\chi^{-1},\phi)$. For $G=SL(2)$ there is one more term, a point mass $c_\tau f^\vee(\chi_0)\tr\tau^\vee(\phi)$ at the unramified quadratic character. Both cases are discussed in \S\S\ref{s:pgl2}--\ref{s:sl2}.
\end{example}

% =====================================================================
\section{Comparison with Route A}\label{s:compare}

This section compares the density $I_B$ of Theorem~\ref{t:B} with the $T$-independent part of the truncated kernel of Route~A. The comparison is Conjecture~\ref{c:AB}. We explain why it should hold, how the two truncations are related, and what Hypothesis (R) of the companion paper means from the present point of view.

\begin{conjecture}\label{c:AB}
For regular unitary $\chi$, the $T$-independent part $I(\chi,\phi)$ of the truncated kernel of \cite[Prop.~4.2]{JC26} equals $I_B(\chi,\phi)$. Equivalently,
\[
I(\chi,\phi)\,\rho(\chi)=c_G\,J_{M_0}(\chi^{-1},\phi).
\]
\end{conjecture}

\paragraph{Why it should hold.} Both sides are linear in $\phi$ and determined by their pairings with all $f^\vee$, and Theorem~4.5(a) of \cite{JC26} gives $I(f,\phi)$ as $\int f^\vee I^T\,d\mu$ for $T$ large. If the oscillatory part of $I^T$ pairs to zero in the limit, which happens for $PGL(2)$ and, numerically, for $PGL(3)$ (\cite[App.~A]{JC26}), then $\int f^\vee I\,d\mu=I(f,\phi)$. Comparing with Theorem~\ref{t:main} for all $f$ then gives Conjecture~\ref{c:AB} when there is no singular part, because the $f^\vee$ are dense in $C(\Thu/W)$. So the conjecture is already a consequence of the PGL(2) results of \cite{JC26} and Theorem~\ref{t:main}, and of the PGL(3) computation, once the normalizations are fixed. The real content is a direct proof that identifies the two constructions.

\paragraph{The two truncations.} Arthur's $u(x,T)$ and the set $G(T)$ of Route~A are both defined by the convex hull of $W\cdot T$. Route~A's vertex terms $\psi(T_Q)R_Q$ \cite[Lemma~3.5A]{JC26} are periodic versions of Arthur's $e^{\Lambda(T_Q)}c_Q(\Lambda)/\theta_Q(\Lambda)$. In rank one, the constant term at $\Lambda=0$ of $\sum_Qe^{\Lambda(T_Q)}c_Q(\Lambda)\theta_Q(\Lambda)^{-1}$ is a logarithmic derivative of an intertwining operator. This is the mechanism that produces a weighted character, and it is visible in \cite[(6.5)]{JC26}.

\paragraph{What (R) means here.} Hypothesis (R) of \cite{JC26} asks that each oscillatory term of $I^T$, times $\rho$, be smooth. Theorem~\ref{t:C} says that only elliptic data on non-trivial $R$-groups can leave a singular part. For $PGL(3)$ all $R$-groups are trivial, and the poles found in \cite[App.~A]{JC26} on the subtori $z_i^3=1$ must cancel in the limit, as observed. The weaker condition (R$'$) of \cite[Rem.~2026.20]{JC26} is a candidate for the analytic shadow of ``no elliptic data''.

\begin{example}[$PGL(2)$]\label{ex:compare}
For $G=PGL(2)$ and $\gamma$ of type (U$_r$), Route~A gives $I(\chi,\gamma)\rho(\chi)$ in closed form. The formula is displayed in \S\ref{s:pgl2}. Conjecture~\ref{c:AB} says that this explicit function equals $-\frac12J_A(\chi^{-1},\phi)$ for $\phi$ concentrated near $\gamma$. By Example~\ref{ex:wchar}, the second function is a trace against a logarithmic derivative of an intertwining operator. The numerical test of this equality is item~\ref{n:5} of \S\ref{s:status}.
\end{example}

% =====================================================================
\section{\texorpdfstring{$SL(n)$}{SL(n)} and endoscopy}\label{s:endo}

This section outlines how, for $G=SL(n)$, the singular part of the spectral measure should be expressed by $\kappa$-orbital integrals and endoscopic groups. It is an outline and not a proof.

For $G=SL(n)$ and $\kappa$ an unramified character of $PGL(n,F)$ of order $m\mid n$, the $\kappa$-orbital integrals and the elliptic endoscopic groups $H_m$ are as in \cite[\S5]{JC26}. Arthur's local character relations \cite{Ar96} and his transfer of weighted orbital integrals \cite{Ar99} relate the elliptic virtual characters $\tau=(M_0,\chi,r)$ of \S\ref{s:class} to the stable characters of the $H_m$: the elliptic element $r\in R_\chi$ corresponds to a $\kappa$ of order equal to that of $r$. Combined with Theorem~\ref{t:main}, this would prove Theorem~\ref{t:D} and remove the hypothesis of \cite[Prop.~5.2]{JC26}. The stabilization of the local trace formula (Arthur's stable trace formula papers; see also Peng's explicit stable local trace formula \cite{Pe16}) is the natural framework.

\begin{status}
This section is an outline. For $\kappa$ of order $n$ the endoscopic group is a torus and the result is \cite[Prop.~5.1]{JC26}.
\end{status}

\begin{example}[$n=2$]\label{ex:endo}
For $G=SL(2)$ there is one non-trivial unramified $\kappa$, of order $2$. It corresponds to the elliptic element of $R_{\chi_0}$ found in Example~\ref{ex:R}. Since $\kappa$ has order $n=2$, the endoscopic group is a torus. The singular part is the point mass at $\chi_0$ of \S\ref{s:sl2}, whose weight $\frac12$ is given by \cite[Prop.~6.4]{JC26}.
\end{example}

% =====================================================================
\section{Examples}\label{s:ex}

This section works out what the results say in four cases. For $PGL(2)$ and $PGL(3)$ there is no singular part, and the density can be compared with the closed forms of Route~A. For $SL(2)$ there is an atom. For $SL(3)$ we state a prediction that has not yet been tested.

\subsection{\texorpdfstring{$PGL(2)$}{PGL(2)}}\label{s:pgl2}

This is the simplest case: rank one, with connected center. Only the absolutely continuous part appears.

Here $Z(G)$ is connected, all $R$-groups are trivial, and Theorem~\ref{t:main} has only the $M_0$ term. Theorem~\ref{t:B} says that $I_B\rho$ is the smooth function $-\frac12J_{A}(\chi^{-1},\phi)$. By \cite[Thm.~6.1, Rem.~6.2]{JC26}, for $\gamma$ of type (U$_r$),
\[
I(\chi,\gamma)\rho(\chi)=\frac Q2\cdot\frac{|B_r|(2-2\cos\theta)+q^r(2\cos\theta-2q^{-1})}{1-2q^{-1}\cos\theta+q^{-2}}\qquad(y=e^{i\theta}),
\]
which is smooth, as Theorem~\ref{t:B} predicts. Conjecture~\ref{c:AB} identifies this with a weighted character.

\begin{status}
Next step: compute $J_A(\chi^{-1},\phi)$ directly, for $\phi$ the normalized characteristic function of $\gamma K_1$. This is a trace on the $(q+1)$-dimensional space of level-one vectors, which can be done numerically, and the result compared with the display above. This will fix the constant $c_G$ and the sign.
\end{status}

\subsection{\texorpdfstring{$SL(2)$}{SL(2)}}\label{s:sl2}

This is the simplest case with a singular part. It comes from the one non-trivial $R$-group of Example~\ref{ex:R}.

At the unramified quadratic character $\chi_0$, $R_{\chi_0}=\{1,s\}$ and $s$ acts by $-1$ on $\fa_0$, so $s$ is elliptic. The corresponding virtual character is $\tau=\pi^+-\pi^-$, where $I(\chi_0)=\pi^+\oplus\pi^-$ and $\pi^+$ contains the spherical vector. So $\tr\tau(f)=f^\vee(\chi_0)$, and the $M=G$ term of Theorem~\ref{t:main} is a point mass at $\chi_0$ with weight $c_\tau\tr\tau^\vee(\phi)$. For $\phi$ concentrated near an unramified elliptic $\gamma$ of depth $0$, $\tr\tau^\vee$ is the $\kappa$-character, and \cite[Prop.~6.4]{JC26} gives the weight $\frac12$. This fixes $c_\tau$ for $SL(2)$, and the result agrees with Route~A.

\subsection{\texorpdfstring{$PGL(3)$}{PGL(3)}}\label{s:pgl3}

This is a rank-two case with connected center. Route~A has a closed form and numerical results to compare with.

All $R$-groups are trivial, so Theorem~\ref{t:C} predicts no singular part, in agreement with the numerical result of \cite[App.~A]{JC26}. Theorem~\ref{t:B} predicts that $I\rho$ is smooth, which agrees with the closed form \cite[(A.4)]{JC26}:
\[
I(\chi,\gamma)\rho(\chi)=\frac{(q-1)^2(q^2+q+1)}{2q^6}\cdot\frac{2(q^2+q+1)-q\sum_{\alpha\in\Phi}z^\alpha}{\prod_{\alpha\in\Phi}(1-q^{-1}z^\alpha)}.
\]
Conjecture~\ref{c:AB} predicts that this is $c_GJ_{M_0}(\chi^{-1},\phi)$ for $\phi$ concentrated near $\gamma$.

\subsection{\texorpdfstring{$SL(3)$}{SL(3)}: a prediction}\label{s:sl3}

This is a rank-two case with non-trivial $R$-groups. Nothing is proved here. We state what Theorem~\ref{t:C} and the outline of \S\ref{s:endo} predict.

Let $G=SL(3)$, and let $\gamma$ be regular elliptic in the norm-one torus of the unramified cubic extension. Up to $W$, the unitary unramified characters with non-trivial $R$-group and an elliptic element in it are the points
\[
\chi_\zeta=(1,\zeta,\zeta^2)\qquad(\zeta^3=1,\ \zeta\ne1),
\]
in coordinates where $\chi$ is a triple modulo scalars (as for the dual group $PGL(3,\mathbb{C})$). At $\chi_\zeta$ a $3$-cycle in $W$ fixes $\chi$ and acts on $\fa_0$ without fixed vectors. Here $R_{\chi_\zeta}\cong\mathbb{Z}/3$, $I(\chi_\zeta)$ has three constituents, and there are two elliptic virtual characters, attached to the two non-trivial elements of $R$. The Satake parameter of $\chi_\zeta$ is $\diag(1,\zeta,\zeta^2)$, which is also the parameter $\chi_\kappa$ attached to the characters $\kappa$ of order $3$ in \cite[Prop.~5.1]{JC26}.

\emph{Prediction.} The spectral measure of $O_\gamma$ on $SL(3)$ is $I\,d\mu$ plus a point mass at $\chi_\zeta$. Its weight is the sum of the two $\kappa$-contributions, $\frac13(\Delta_\kappa^{-1}+\Delta_{\kappa^2}^{-1})c_S$ in the notation of \cite[(5.4A), Prop.~5.1]{JC26}, which is $\frac23$ if all transfer factors and $c_S$ equal $1$. There are no singular contributions on one-dimensional subtori: those would require an elliptic element of an $R$-group for $M\cong GL(2)$, and the $R$-groups of unramified principal series of $GL(2)$ are trivial.

\begin{status}
To be tested numerically on the building of $SL(3)$, by adapting the $PGL(3)$ code of \cite[App.~A]{JC26} (\texttt{examples/pgl3/}) to vertices of one type and to the cubic norm-one torus. The predicted mass $\frac23$ depends on normalizations that have not been checked.
\end{status}

% =====================================================================
\section{Summary of results}\label{s:summary}

This section collects the results of the paper, with the status of each as it is stated where the result appears. No status here is meant to be stronger than in the body. The remaining work and the open questions are listed in \S\ref{s:status}.

\begin{center}\small
\begin{tabular}{@{}>{\raggedright\arraybackslash}p{0.21\textwidth}>{\raggedright\arraybackslash}p{0.33\textwidth}>{\raggedright\arraybackslash}p{0.39\textwidth}@{}}
\toprule
Result & Content & Status\\
\midrule
Lemma~\ref{l:sphvanish} & Weighted characters of spherical functions vanish for $M\ne L$. & Proved here from (\ref{e:fix}). The comparison of the normalization with Arthur's conditions is item~\ref{n:1}.\\[0.5ex]
Proposition~\ref{p:geom} & The geometric side reduces to $I(f,\phi)$. & Proved here. The constants in (\ref{e:geo}) are being checked. See item~\ref{n:2}.\\[0.5ex]
Proposition~\ref{p:split} & The spherical factor contributes only $f^\vee(\chi)$. & Proved here from (\ref{e:split}), whose form in the local setting is being checked. See item~\ref{n:2}.\\[0.5ex]
Lemma~\ref{l:cusp} & Functions supported on $G_{\ellr}$ are cuspidal. & Proved here from van Dijk's formula \cite{vD72}.\\[0.5ex]
Lemma~\ref{l:noDS} & Square-integrable representations of Levi subgroups $L\ne M_0$ have no hyperspecial-fixed vectors. & Proof sketched. A reference or a complete proof is to be supplied. See item~\ref{n:3}.\\[0.5ex]
Theorem~\ref{t:main} and Theorems~\ref{t:A}, \ref{t:B}, \ref{t:C} & The spectral measure exists. Its absolutely continuous part is a weighted character, and its singular part comes from elliptic virtual characters. & Proof in outline, with the normalizations of \S\ref{s:status} assumed. See items~\ref{n:1}, \ref{n:2} and \ref{n:4}.\\[0.5ex]
Theorem~\ref{t:D} & For $SL(n)$ the singular part is given by endoscopy. & Outline only. See \S\ref{s:endo} and item~\ref{n:8}.\\[0.5ex]
Conjecture~\ref{c:AB} & The density $I_B$ equals the $T$-independent part of the kernel of Route~A. & Conjectural. See item~\ref{n:6}.\\[0.5ex]
\S\ref{s:sl3} & A point mass at $\chi_\zeta$ for $SL(3)$. & A prediction. It has not been tested numerically, and it depends on normalizations that have not been checked. See item~\ref{n:7}.\\
\bottomrule
\end{tabular}
\end{center}

The examples of \S\ref{s:ex} for $PGL(2)$, $SL(2)$ and $PGL(3)$ compare Theorem~\ref{t:main} with results of \cite{JC26}. The rank-one computation that would fix the constants $c_G$ and $c_\tau$ directly is item~\ref{n:5}.

% =====================================================================
\section{Status and what remains}\label{s:status}

This section lists the steps that remain before the results of the paper are fully proved. The items are cited by number throughout the paper. They also serve as the list of open questions.

\begin{enumerate}[label=(N\arabic*)]
\item\label{n:1} \emph{Normalizing factors.} Check that (\ref{e:rnorm}) satisfies Arthur's conditions for normalizing factors and agrees with \cite{Ar98} in the unramified case (non-split groups included).
\item\label{n:2} \emph{Arthur's formula.} Fix the precise statements, constants and measures of the geometric and spectral expansions and of the elliptic form from \cite{Ar91,Ar93}, and the splitting formula (\ref{e:split}) in the local setting.
\item\label{n:3} \emph{Lemma~\ref{l:noDS}.} Supply a reference or a complete proof.
\item\label{n:4} \emph{Theorem~\ref{t:main}.} Write the proof in full, including the smoothness of $\chi\mapsto J_{M_0}(\chi,\phi)$ in the compact picture.
\item\label{n:5} \emph{Rank one.} Compute $J_A(\chi,\phi)$ for $PGL(2)$ and $SL(2)$ numerically, and compare with \cite[Thm.~6.1, Prop.~6.4]{JC26}. This fixes $c_G$ and $c_\tau$.
\item\label{n:6} \emph{Conjecture~\ref{c:AB}.} Give a direct proof, at least in rank one.
\item\label{n:7} \emph{$SL(3)$.} Run the numerical test of the prediction.
\item\label{n:8} \emph{Theorem~\ref{t:D}.} Write \S\ref{s:endo} in full.
\end{enumerate}

% =====================================================================
\appendix
\section{Glossary}\label{s:gloss}

This glossary gives short definitions of the terms used in this paper, with pointers to where they are used. Terms already in the glossary of \cite{JC26} are repeated only briefly.

\begin{description}[leftmargin=1em,style=nextline,itemsep=3pt]
\item[Unramified group; hyperspecial $K$] A reductive group over $F$, quasi-split and split over an unramified extension. $K$ is the stabilizer of a hyperspecial vertex of the building; for split $G$, $K=G(\OO)$.
\item[Spherical function $\Gamma_\chi$; Satake transform $f^\vee$] $\Gamma_\chi$ is the bi-$K$-invariant matrix coefficient of the spherical vectors of $I_B(\chi)$. The Satake transform is $f^\vee(\chi)=\tr I_B(\chi,f)$ for $f\in\Hc(G,K)$, a $W$-invariant Laurent polynomial in $\chi$.
\item[Plancherel measure $d\mu=\rho\,d\chi$] The measure on $\Thu/W$ for which $f\mapsto f^\vee$ is an isometry from $L^2(K\backslash G/K)$ (Macdonald). $\rho$ is proportional to $|c(\chi)|^{-2}$.
\item[Principal series $I_P(\chi)$; compact picture] Normalized induction from a character of $M_0(F)$, realized on functions on $K$ so that all $I_P(\chi)$ share one space.
\item[Intertwining operator $J_{P'|P}(\chi)$] The operator $I_P(\chi)\to I_{P'}(\chi)$ given by integration over $N_{P'}\cap N_{\bar P}$, meromorphic in $\chi$ (\S\ref{s:norm}).
\item[Normalizing factor; normalized operator $R_{P'|P}$] A scalar $r_{P'|P}(\chi)$ with $R=r^{-1}J$ unitary and holomorphic on the unitary axis and multiplicative. In the unramified case it is the product of $L(0,\chi_\alpha)/L(1,\chi_\alpha)$ over the relevant roots, and $R$ fixes the spherical vector (\ref{e:fix}).
\item[$(G,M)$-family] A family $\{c_Q(\Lambda)\}_{Q\in\cP(M)}$ of smooth functions that agree on common walls. Its limit $c_M=\lim_{\Lambda\to0}\sum_Qc_Q(\Lambda)\theta_Q(\Lambda)^{-1}$ exists. For a constant family $c_M=0$ unless $M=G$ \cite{Ar81}.
\item[Weighted orbital integral $J_M(\gamma,f)$] An orbital integral against the weight function $v_M$ attached to a $(G,M)$-family. It is not invariant unless $M=G$. These are the terms of the geometric side of Arthur's trace formulas.
\item[Weighted character $J_M(\pi,f)$] $\tr(\mathcal R_M(\pi,P)I_P(\pi,f))$, with $\mathcal R_M$ the limit of the $(G,M)$-family of normalized intertwining operators (\S\ref{s:wchar}). These are the terms of the spectral side. They depend on the normalization.
\item[Splitting formula] The expression (\ref{e:split}) of a weighted character of $f_1\otimes f_2$ through weighted characters of $f_1$ and of $f_2$ on Levi subgroups, with coefficients $d_M^G(L_1,L_2)$.
\item[Local trace formula] Arthur's identity between the geometric expansion (\ref{e:geo}) and the spectral expansion (\ref{e:spec}) of the constant term $J(f)$ of a truncated integral of the kernel of $R(f)$ on $L^2(G(F))$, for $f\in C_c^\infty(G(F)\times G(F))$ \cite{Ar91}.
\item[Truncation $u(x,T)$] The characteristic function of the compact set of $x\in G(F)$ whose Cartan component lies in the convex hull of $W\cdot T$. It is the analogue of the set $G(T)$ of Route~A.
\item[Regular elliptic element; $G_{\ellr}$] A regular semisimple element whose centralizer is an anisotropic torus (modulo the center).
\item[Cuspidal function] $\phi\in C_c^\infty(G(F))$ whose orbital integrals vanish on non-elliptic regular elements; equivalently, $\tr\sigma(\phi)=0$ for properly induced $\sigma$. Functions supported on $G_{\ellr}$ are cuspidal (Lemma~\ref{l:cusp}).
\item[$R$-group $R_\chi$] $W_\chi/W^\circ_\chi$, the stabilizer of $\chi$ in $W$ modulo the subgroup generated by the reflections in roots $\alpha$ with $\chi_\alpha=1$. It controls the reducibility of $I_P(\chi)$ (\S\ref{s:R}).
\item[Elliptic element of an $R$-group] An $r\in R_\chi$ represented by $w$ with no non-zero fixed vector in $\fa_0^G$ (or in $\fa_0^M$, for $M$).
\item[$\Tdisc(M)$, $\Tell(M)$] Arthur's sets of virtual tempered characters: triples $(L,\sigma,r)$ with $\sigma$ square-integrable on a Levi $L\subseteq M$ and $r$ in the $R$-group of $\sigma$ in $M$. $\Tell(M)$ is the subset with $r$ elliptic \cite{Ar93}.
\item[Absolutely continuous and singular parts] The decomposition $dI_\phi=I_B\,d\mu+dI^{\mathrm{sing}}_\phi$ of the spectral measure (Theorems~\ref{t:B} and \ref{t:C}).
\item[$\kappa$-orbital integral; endoscopic group; transfer factor] See \cite[\S5]{JC26}. For $SL(n)$, $\kappa$ is a character of $PGL(n,F)$ trivial on the image of $SL(n,F)$. The elliptic endoscopic groups are the $H_m$, and transfer factors $\Delta$ relate their stable orbital integrals to $\kappa$-orbital integrals.
\item[Fundamental lemma] The identity, for the spherical Hecke algebra, between $\kappa$-orbital integrals and stable orbital integrals on endoscopic groups. It is a theorem (Ng\^o \cite{Ng10} and others).
\item[Route A, Route B] Route A is the truncation method of \cite{JC26}. Route B, this paper, specializes Arthur's local trace formula.
\item[Hypothesis (R)] The regularity hypothesis of \cite[\S4.3]{JC26}: each oscillatory coefficient of the truncated kernel, times $\rho$, is smooth. It is sufficient for the spectral identity but not necessary (\S\ref{s:compare}).
\end{description}

\section*{Acknowledgments}
\addcontentsline{toc}{section}{Acknowledgments}
This paper began with a suggestion of Claude (Anthropic), an AI model. During the 2026 revision of the first author's paper \cite{Jo92}, Claude suggested reproving the results of the 1997 version of that paper from Arthur's later work on the local trace formula \cite{Ar91,Ar93}. The first author decided to pursue the suggestion, directed the work and reviewed the drafts.

Claude took part as a co-author. It wrote the outline (\texttt{OUTLINE.md}, in the directory of the source), which the plan of the paper follows. It drafted the text, the examples, the literature review, the summary and the glossary, and it made the style revision of Draft~2. No other AI system contributed, and the paper has not yet been refereed.

Review by AI systems is not an independent certification of the results. The errors that remain are the authors'. Further details on the AI usage, and a list of every change between drafts, can be found in the revision log (\texttt{revision\_log.md}) that accompanies the source.


\begin{thebibliography}{BDK86}
\bibitem[Ar81]{Ar81} J.~Arthur, \emph{The trace formula in invariant form}, Ann.\ of Math.\ \textbf{114} (1981), 1--74.
\bibitem[Ar88]{Ar88} J.~Arthur, \emph{The invariant trace formula I. Local theory}, J.\ Amer.\ Math.\ Soc.\ \textbf{1} (1988), 323--383.
\bibitem[Ar89a]{Ar89a} J.~Arthur, \emph{Towards a local trace formula}, in: Algebraic Analysis, Geometry, and Number Theory, Johns Hopkins Univ.\ Press, Baltimore, 1989, 1--24.
\bibitem[Ar89b]{Ar89b} J.~Arthur, \emph{Intertwining operators and residues I. Weighted characters}, J.\ Funct.\ Anal.\ \textbf{84} (1989), 19--84.
\bibitem[Ar91]{Ar91} J.~Arthur, \emph{A local trace formula}, Publ.\ Math.\ IH\'ES \textbf{73} (1991), 5--96.
\bibitem[Ar93]{Ar93} J.~Arthur, \emph{On elliptic tempered characters}, Acta Math.\ \textbf{171} (1993), 73--138.
\bibitem[Ar94]{Ar94} J.~Arthur, \emph{On the Fourier transforms of weighted orbital integrals}, J.\ reine angew.\ Math.\ \textbf{452} (1994), 163--217.
\bibitem[Ar96]{Ar96} J.~Arthur, \emph{On local character relations}, Selecta Math.\ (N.S.) \textbf{2} (1996), 501--579.
\bibitem[Ar98]{Ar98} J.~Arthur, \emph{Canonical normalization of weighted characters and a transfer conjecture}, C.\ R.\ Math.\ Rep.\ Acad.\ Sci.\ Canada \textbf{20} (1998), 35--52.
\bibitem[Ar99]{Ar99} J.~Arthur, \emph{On the transfer of distributions: weighted orbital integrals}, Duke Math.\ J.\ \textbf{99} (1999), 209--283.
\bibitem[Ar01]{Ar01} J.~Arthur, \emph{A stable trace formula II. Global descent}, Invent.\ Math.\ \textbf{143} (2001), 157--220.
\bibitem[Ar02]{Ar02} J.~Arthur, \emph{A stable trace formula I. General expansions}, J.\ Inst.\ Math.\ Jussieu \textbf{1} (2002), 175--277.
\bibitem[Ar03]{Ar03} J.~Arthur, \emph{A stable trace formula III. Proof of the main theorems}, Ann.\ of Math.\ \textbf{158} (2003), 769--873.
\bibitem[Ar13]{Ar13} J.~Arthur, \emph{The Endoscopic Classification of Representations: Orthogonal and Symplectic Groups}, Amer.\ Math.\ Soc.\ Colloquium Publications \textbf{61}, 2013.
\bibitem[BDK86]{BDK86} J.~Bernstein, P.~Deligne, D.~Kazhdan, \emph{Trace Paley--Wiener theorem for reductive $p$-adic groups}, J.\ Analyse Math.\ \textbf{47} (1986), 180--192.
\bibitem[BP20]{BP20} R.~Beuzart-Plessis, \emph{A local trace formula for the Gan--Gross--Prasad conjecture for unitary groups: the archimedean case}, Ast\'erisque \textbf{418} (2020).
\bibitem[BZ77]{BZ77} I.~N.~Bernstein, A.~V.~Zelevinsky, \emph{Induced representations of reductive $p$-adic groups.\ I}, Ann.\ Sci.\ \'Ecole Norm.\ Sup.\ (4) \textbf{10} (1977), 441--472.
\bibitem[Ca79]{Ca79} P.~Cartier, \emph{Representations of $p$-adic groups: a survey}, in: Automorphic Forms, Representations and $L$-functions (Corvallis), Proc.\ Sympos.\ Pure Math.\ \textbf{33}, Part~1, Amer.\ Math.\ Soc., 1979, 111--155.
\bibitem[Ca80]{Ca80} W.~Casselman, \emph{The unramified principal series of $\mathfrak p$-adic groups I. The spherical function}, Compositio Math.\ \textbf{40} (1980), 387--406.
\bibitem[Ca95]{Ca95} W.~Casselman, \emph{Introduction to the theory of admissible representations of $p$-adic reductive groups}, unpublished notes, 1974; revised draft 1995.
\bibitem[vD72]{vD72} G.~van Dijk, \emph{Computation of certain induced characters of $\mathfrak p$-adic groups}, Math.\ Ann.\ \textbf{199} (1972), 229--240.
\bibitem[JC26]{JC26} D.~Joyner and Claude, \emph{Simple local trace formulas for unramified $p$-adic groups. Draft 2. Corrected and expanded version}, preprint, November 2026 (a revision of \cite{Jo92} by the truncation method; ``Route~A''), \url{https://github.com/wdjoyner/trace-formulas/tree/main/local-trace/paper}.
\bibitem[Jo92]{Jo92} D.~Joyner, \emph{Simple local trace formulas for unramified $p$-adic groups}, Pacific J.\ Math.\ \textbf{154} (1992), 103--129.
\bibitem[Ka86]{Ka86} D.~Kazhdan, \emph{Cuspidal geometry of $p$-adic groups}, J.\ Analyse Math.\ \textbf{47} (1986), 1--36.
\bibitem[KS71]{KS71} A.~W.~Knapp, E.~M.~Stein, \emph{Intertwining operators for semisimple groups}, Ann.\ of Math.\ \textbf{93} (1971), 489--578.
\bibitem[Ke82]{Ke82} C.~D.~Keys, \emph{On the decomposition of reducible principal series representations of $p$-adic Chevalley groups}, Pacific J.\ Math.\ \textbf{101} (1982), 351--388.
\bibitem[La76]{La76} R.~P.~Langlands, \emph{On the Functional Equations Satisfied by Eisenstein Series}, Lecture Notes in Math.\ \textbf{544}, Springer, 1976.
\bibitem[LL79]{LL79} J.-P.~Labesse, R.~P.~Langlands, \emph{$L$-indistinguishability for $SL(2)$}, Canad.\ J.\ Math.\ \textbf{31} (1979), 726--785.
\bibitem[LS87]{LS87} R.~P.~Langlands, D.~Shelstad, \emph{On the definition of transfer factors}, Math.\ Ann.\ \textbf{278} (1987), 219--271.
\bibitem[Li21]{Li21} H.~Li, \emph{A local trace formula for $p$-adic infinitesimal symmetric spaces: the case of Guo--Jacquet}, preprint, arXiv:2105.00087.
\bibitem[Ma71]{Ma71} I.~G.~Macdonald, \emph{Spherical Functions on a Group of $p$-adic Type}, Publ.\ Ramanujan Inst.\ \textbf{2}, Univ.\ of Madras, 1971.
\bibitem[MW16]{MW16} C.~M\oe glin, J.-L.~Waldspurger, \emph{Stabilisation de la formule des traces tordue}, Vols.~1--2, Progress in Math.\ \textbf{316}--\textbf{317}, Birkh\"auser, 2016.
\bibitem[Ng10]{Ng10} B.~C.~Ng\^o, \emph{Le lemme fondamental pour les alg\`ebres de Lie}, Publ.\ Math.\ IH\'ES \textbf{111} (2010), 1--169.
\bibitem[OS13]{OS13} E.~Opdam, M.~Solleveld, \emph{Extensions of tempered representations}, Geom.\ Funct.\ Anal.\ \textbf{23} (2013), 664--714.
\bibitem[Pe16]{Pe16} Z.~Peng, \emph{Multiplicity formula and stable trace formula}, preprint, arXiv:1608.00055.
\bibitem[Re01]{Re01} M.~Reeder, \emph{Euler--Poincar\'e pairings and elliptic representations of Weyl groups and $p$-adic groups}, Compositio Math.\ \textbf{129} (2001), 149--182.
\bibitem[Sa63]{Sa63} I.~Satake, \emph{Theory of spherical functions on reductive algebraic groups over $\mathfrak p$-adic fields}, Publ.\ Math.\ IH\'ES \textbf{18} (1963), 5--69.
\bibitem[Si78]{Si78} A.~J.~Silberger, \emph{The Knapp--Stein dimension theorem for $p$-adic groups}, Proc.\ Amer.\ Math.\ Soc.\ \textbf{68} (1978), 243--246.
\bibitem[Wa89a]{Wa89a} J.-L.~Waldspurger, \emph{Int\'egrales orbitales sph\'eriques pour $GL(N)$ sur un corps $p$-adique}, Ast\'erisque \textbf{171--172} (1989), 279--337.
\bibitem[Wa03]{Wa03} J.-L.~Waldspurger, \emph{La formule de Plancherel pour les groupes $p$-adiques (d'apr\`es Harish-Chandra)}, J.\ Inst.\ Math.\ Jussieu \textbf{2} (2003), 235--333.
\bibitem[Wa10]{Wa10} J.-L.~Waldspurger, \emph{Une formule int\'egrale reli\'ee \`a la conjecture locale de Gross--Prasad}, Compositio Math.\ \textbf{146} (2010), 1180--1290.
\end{thebibliography}

\end{document}
